\section{Discussion and Future Work}

% The original Platonic Representation Hypothesis conjectures that the representation spaces of modern neural networks are converging. While the hypothesis itself is up for debate, it is even less obvious that there exists a method to explicitly discover such a universal representation. Remarkably, we demonstrate not only that this universal representation exists (across the encoders we study), but that it can be systematically constructed using our method.


% This robust in-distribution and out-of-distribution performance indicates that \atk{} does not just learn to translate its training embeddings but rather learns a domain‑agnostic mapping. By encoding the underlying semantics of models $M_1$ and $M_2$ into its shared latent space, our method learns a mathematical relationship between embedding spaces that transfers across diverse data distributions. This allows \atk{} to generate the correct embedding for input text it has never seen on an encoder it does not have access to, just by harnessing the universal geometry of embeddings.

% \atk{}’s strong performance on both in- and out-of-distribution data suggests it does more than memorize training translations---it learns the mathematical relationship between embedding spaces. By capturing the shared semantics of models $M_1$ and $M_2$ in a universal latent space, \atk{} can generate accurate target embeddings for unseen text from an encoder it has never encountered, purely by leveraging the universal geometry of embeddings.

% Importantly, we were able to extract sensitive attributes (disease information) from patient records \textit{with access to only their raw embeddings}, implying that a leak of hospital embeddings is equivalent to leakage of patient data. As translation methods like \atk{} improve, higher-fidelity extractions (perhaps of names, addresses, and other personal data) will become easier and more harmful. This trend necessitates that practitioners start to handle embeddings like they would handle their underlying text.

% The quality of inversions will improve dramatically with improvements in both translation and inversion. As such, we again emphasize the importance of treating embeddings safely---less like the raw numbers they appear as and more like the underlying texts they are functionally equivalent to.

% We emphasize that our method measures a \textit{lower bound} on the potential of inter-embedding translation, and any methodological improvements (including both better learning algorithms and better architectures) will improve our ability to translate between modalities. In particular, a class of more stable models would enable us to scale our system to more parameters and more data, likely improving our ability to translate in most settings.

% Our experiments relate to text-only embeddings as well as multimodal models like CLIP that include text as a modality. We see our results as empirical validation of the Platonic Representation Hypothesis \citep{huh2024platonicrepresentationhypothesis}, and look forward to other future work that extends our framework to new modalities and model families.

The Platonic Representation Hypothesis conjectures that the representation spaces of modern neural networks are converging.   We assert the Strong Platonic Representation Hypothesis: the latent universal representation can be learned and harnessed to translate between representation spaces without any encoders or paired data.

In \Cref{sec:eval_geometry}, we demonstrated 
that our \atk{} method successfully translates embeddings generated from unseen documents by unseen encoders, and the translator is robust to (sometimes very) out-of-distribution inputs.  This suggests that \atk{} learns domain-agnostic translations based on the universal geometric relationships which encode the same semantics in multiple embedding spaces.

% This suggests our method doesn't just memorize training examples but learns a domain-agnostic mapping that captures the underlying mathematical relationship between embedding spaces. By encoding the shared semantics of models $M_1$ and $M_2$ into a universal latent space, \atk{} can generate accurate target embeddings for unseen text from unfamiliar encoders by harnessing the universal geometry of embeddings.

In \Cref{sec:eval_semantics}, we showed that \atk{} translations preserve sufficient input semantics to enable attribute inference. 
We extracted sensitive disease information from patient records and partial content from corporate emails, with access only to document embeddings and no access to the encoder that produced them.  Better translation methods will enable higher-fidelity extraction, confirming once again that embeddings reveal (almost) as much as their inputs.

% We successfully extracted sensitive disease information from patient records and confidential information from business records with access only to their raw embeddings. This demonstrates that embedding leakage is functionally equivalent to leaking the underlying text itself. As translation methods like \atk{} improve, higher-fidelity extractions of personal data will become increasingly feasible, underscoring the importance of treating embeddings as sensitive as the texts they represent.

Our findings provide compelling evidence for the Strong Platonic Representation Hypothesis for text-based models.  Our preliminary results on CLIP suggest that the universal geometry can be harnessed in other modalities, too.  The results in this paper are but a \emph{lower bound} on inter-representation translation.  Better and more stable learning algorithms, architectures, and other methodological improvements will support scaling to more data, more model families, and more modalities.

% Based on our findings, we see our results as compelling evidence for both the original and Strong Platonic Representation Hypotheses for text-based models. Our preliminary results on CLIP further extend this perspective, leading us to conjecture that the full Strong Platonic Representation Hypothesis holds true. Importantly, our method demonstrates a \textit{lower bound} on the potential of inter-embedding translation, and any methodological improvements (including both better learning algorithms and architectures) will improve our ability to translate between modalities. In particular, a class of more stable models would enable us to scale our system to more parameters and more data, likely improving our ability to translate in most settings. We look forward to future work that extends our framework to new modalities and model families.